\begin{table}[H]
\centering
\captionsetup{justification=centering}
\caption{Hiperparâmetros estimados para a taxa de desocupação - modelos univariado e multivariado}
\label{tab:hipertaxa}
\scalebox{0.92}{
\renewcommand{\arraystretch}{0.85}
\begin{tabular}{llccccc}
\toprule
& & \multicolumn{5}{c}{\textbf{Modelo Univariado}} \\
\cmidrule(lr){3-7}
\textbf{Estrato geográfico} & \textbf{Processo} & \(\hat{\sigma}_L^2\) & \(\hat{\sigma}_R^2\) & \(\hat{\sigma}_S^2\) & \(\hat{\sigma}_I^2\) & \(\hat{\sigma}_{\tilde{e}}^2\) \\
\midrule
01 - Belo Horizonte & AR(1) & 0,0000 & 0,1028 & 0,0000 & 0,0000 & 0,9263 \\
02 - Entorno e Colar Metropolitano de BH & MA(1) & 0,0000 & 0,3350 & 0,0000 & 0,0000 & 0,9485 \\
03 - Sul de Minas & AR(1) & 0,5334 & 0,0051 & 0,0000 & 0,0000 & 0,8744 \\
04 - Triângulo Mineiro & AR(1) & 0,0000 & 0,0431 & 0,0000 & 0,0000 & 0,9604 \\
05 - Zona da Mata & AR(1) & 0,3514 & 0,0023 & 0,0000 & 0,0000 & 0,9584 \\
06 - Norte de Minas & MA(1) & 0,1873 & 0,1890 & 0,0091 & 0,0000 & 0,8083 \\
07 - Vale do Rio Doce & AR(1) & 1,0321 & 0,0100 & 0,0000 & 0,0000 & 0,8583 \\
08 - Central & AR(3) & 0,0003 & 0,1653 & 0,0000 & 0,0000 & 0,8978 \\
\midrule
& & \multicolumn{5}{c}{\textbf{Modelo Multivariado}} \\
\cmidrule(lr){3-7}
\textbf{Estrato geográfico} & \textbf{Processo} & \(\hat{\sigma}_L^2\) & \(\hat{\sigma}_R^2\) & \(\hat{\sigma}_S^2\) & \(\hat{\sigma}_I^2\) & \(\hat{\sigma}_{\tilde{e}}^2\) \\
\midrule
01 - Belo Horizonte & AR(1) & 0,0000 & 0,2211 & 0,0000 & 0,0000 & 0,9263 \\
02 - Entorno e Colar Metropolitano de BH & MA(1) & 0,0000 & 0,3694 & 0,0000 & 0,0000 & 0,9485 \\
03 - Sul de Minas & AR(1) & 0,0004 & 0,1506 & 0,0000 & 0,0000 & 0,8744 \\
04 - Triângulo Mineiro & AR(1) & 0,0000 & 0,0876 & 0,0000 & 0,0000 & 0,9604 \\
05 - Zona da Mata & AR(1) & 0,0055 & 0,0814 & 0,0000 & 0,0000 & 0,9584 \\
06 - Norte de Minas & MA(1) & 0,0017 & 0,2484 & 0,0099 & 0,0000 & 0,8083 \\
07 - Vale do Rio Doce & AR(1) & 0,0016 & 0,3222 & 0,0000 & 0,0000 & 0,8583 \\
08 - Central & AR(3) & 0,0003 & 0,1245 & 0,0000 & 0,0000 & 0,8978 \\
\bottomrule
\end{tabular}}
\fonte{Elaboração própria, com base nos dados da PNAD Contínua. Nota: o processo do erro amostral é identificado individualmente por estrato (Seção \ref{sec:metodologia}), e \(\hat{\sigma}_{\tilde{e}}^2\) é derivada dos coeficientes do processo, por isso coincidindo nos dois modelos. Sobre \(\sigma_I^2\), ver a ressalva de identificação no texto.}
\end{table}
